\documentclass[UTF8,11pt]{ctexart}
\usepackage[a4paper,margin=22mm]{geometry}
\usepackage{graphicx,enumitem,amssymb,tabularx}
\usepackage[normalem]{ulem}
\usepackage{xeCJKfntef}
\setlength{\parindent}{0pt}
\setlength{\parskip}{.6em}
\sloppy
\begin{document}
\section*{2016年6月英语六级真题(第3套)}
\subsection*{\textbf{2016} 年\textbf{6} 月大学英语六级考试真题（第三套）}


\subsection*{\textbf{Part I Writing (30 minutes)}}


\textbf{Directions:} For this part, you are allowed 30 minutes to write a short essay on the use of robots. Try to imagine what will happen when more and more robots take the place of human beings in industry as well as people’s daily lives. You are required to write at least 150 words but no more than 200 words.


\subsection*{\textbf{Section A}}


\textbf{Directions:} In this section, you will hear two long conversations. At the end of each conversation, you will hear four questions. Both the conversation and the questions will be spoken only once. After you hear a question, you must choose the best answer. from the four choices marked A), B),C) and D). Then mark the corresponding letter on Answer Sheet 1 with a single line through the centre.


\textbf{Questions 1 to 4 are based on the conversation you have just heard.}


1.


\begin{itemize}[leftmargin=2em,itemsep=.2em,topsep=.4em]\item[A.] It is advertising electronic products.
\item[B.] It is planning to tour East Asia.
\item[C.] It is sponsoring a TV programme.
\item[D.] It is giving performances in town.
\end{itemize}

\noindent\begin{tabularx}{\linewidth}{@{}lXXXX@{}}2. & \textbf{A.} 20,000 pounds. & \textbf{B.} 12,000 pounds. & \textbf{C.} Less than 20,000 pounds. & \textbf{D.} Less than 12,000 pounds.\\\end{tabularx}\par

3.


\begin{itemize}[leftmargin=2em,itemsep=.2em,topsep=.4em]\item[A.] A lot of good publicity.
\item[B.] Talented artists to work for it.
\item[C.] Long-term investments.
\item[D.] A decrease in production costs.
\end{itemize}

4.


\begin{itemize}[leftmargin=2em,itemsep=.2em,topsep=.4em]\item[A.] Promise long-term cooperation with the Company.
\item[B.] Explain frankly their own current financial situation.
\item[C.] Pay for the printing of the performance programme.
\item[D.] Bear the cost of publicising the Company's performance.
\end{itemize}

\textbf{Questions 5 to 8 are based on the conversation you have just heard.}


5.


\begin{itemize}[leftmargin=2em,itemsep=.2em,topsep=.4em]\item[A.] He has been seeing doctors and counsellors.
\item[B.] He has found a new way to train his voice.
\item[C.] He was caught abusing drugs.
\item[D.] He might give up concert tours.
\end{itemize}

6.


\begin{itemize}[leftmargin=2em,itemsep=.2em,topsep=.4em]\item[A.] Singers may become addicted to it.
\item[B.] It helps singers warm themselves up.
\item[C.] Singers use it to stay away from colds.
\item[D.] It can do harm to singers' vocal chords.
\end{itemize}

7.


\begin{itemize}[leftmargin=2em,itemsep=.2em,topsep=.4em]\item[A.] They are eager to become famous.
\item[B.] Many lack professional training.
\item[C.] Few will become successful.
\item[D.] They live a glamorous life.
\end{itemize}

8.


\begin{itemize}[leftmargin=2em,itemsep=.2em,topsep=.4em]\item[A.] Harm to singers done by smoky atmospheres.
\item[B.] Side effects of some common drugs.
\item[C.] Voice problems among pop singers.
\item[D.] Hardships experienced by many young singers.
\end{itemize}

\subsection*{\textbf{Section B}}


\textbf{Directions:} In this section, you will hear two passages. At the end of each passage, you will hear three or four questions. Both the passage and the questions will be spoken only once. After you hear a question, you must choose the best answer from the four choices marked A), B), C) and D). Then mark the corresponding letter on Answer Sheet


1 with a single line through the centre.


\subsection*{\textbf{Passage One}}


\textbf{Questions 9 to 12 are based on the passage you have just heard.}


9.


\begin{itemize}[leftmargin=2em,itemsep=.2em,topsep=.4em]\item[A.] It has not been very successful.
\item[B.] It has long become a new trend.
\item[C.] It has met with strong resistance.
\item[D.] It has attracted a lot of users.
\end{itemize}

10.


\begin{itemize}[leftmargin=2em,itemsep=.2em,topsep=.4em]\item[A.] It saves time.
\item[B.] It increases parking capacity.
\item[C.] It ensures drivers' safety.
\item[D.] It reduces car damage.
\end{itemize}

11.


\begin{itemize}[leftmargin=2em,itemsep=.2em,topsep=.4em]\item[A.] Collect money and help new users.
\item[B.] Maintain the automated system.
\item[C.] Stay alert to any emergency.
\item[D.] Walk around and guard against car theft.
\end{itemize}

12.


\begin{itemize}[leftmargin=2em,itemsep=.2em,topsep=.4em]\item[A.] They will vary with the size of vehicles.
\item[B.] They will be discountable to regular customers.
\item[C.] They will be lower than conventional parking.
\item[D.] They will be reduced if paid in cash.
\end{itemize}

\subsection*{\textbf{Passage Two}}


\textbf{Questions 13 to 15 are based on the passage you have just heard.}


13.


\begin{itemize}[leftmargin=2em,itemsep=.2em,topsep=.4em]\item[A.] They do not know any solution.
\item[B.] They do not give up drunk driving.
\item[C.] They do not behave in public places.
\item[D.] They do not admit being alcohol addicts.
\end{itemize}

14.


\begin{itemize}[leftmargin=2em,itemsep=.2em,topsep=.4em]\item[A.] To stop them from fighting back.
\item[B.] To thank them for their hospitality.
\item[C.] To teach them the European lifestyle.
\item[D.] To relieve their pains and sufferings.
\end{itemize}

15.


\begin{itemize}[leftmargin=2em,itemsep=.2em,topsep=.4em]\item[A.] Without intervention they will be a headache to the nation.
\item[B.] With support they can be brought back to a normal life.
\item[C.] They readily respond to medical treatment.
\item[D.] They pose a serious threat to social stability.
\end{itemize}

\subsection*{\textbf{Section C}}


\textbf{Directions:} In this section, you will hear three recordings of lectures or talks followed by three or four questions. The recordings will be played only once. After you hear a question, you must choose the best answer from the four choices marked A),B),C) and D). Then mark the corresponding letter on Answer Sheet 1 with a single line through the centre.


\textbf{Recording One}


\textbf{Questions 16 to 19 are based on the recording you have just heard.}


16.


\begin{itemize}[leftmargin=2em,itemsep=.2em,topsep=.4em]\item[A.] To award them for their hard work.
\item[B.] To build common views.
\item[C.] To bring in business projects.
\item[D.] To vote for action.
\end{itemize}

17.


\begin{itemize}[leftmargin=2em,itemsep=.2em,topsep=.4em]\item[A.] Recovering from the Great Recession.
\item[B.] Creating jobs and boosting the economy.
\item[C.] Rewarding innovative businesses.
\item[D.] Launching economic campaigns.
\end{itemize}

18.


\begin{itemize}[leftmargin=2em,itemsep=.2em,topsep=.4em]\item[A.] Talking over paying off deficit.
\item[B.] Increasing the number of middle class.
\item[C.] Controlling the impact on education.
\item[D.] Planning to reduce energy consumption.
\end{itemize}

19.


\begin{itemize}[leftmargin=2em,itemsep=.2em,topsep=.4em]\item[A.] Shorten America's way to prosperity.
\item[B.] Be cautious about reducing the deficit.
\item[C.] Increase deficit to cover the revenue.
\item[D.] Require the richest to pay more taxes.
\end{itemize}

\textbf{Questions 20 to 22 are based on the recording you have just heard}.


20.


\begin{itemize}[leftmargin=2em,itemsep=.2em,topsep=.4em]\item[A.] They can be redeemed for cash.
\item[B.] They can be used to reduce meal costs.
\item[C.] They can be used as membership certificate.
\item[D.] They can be used to make reservations.
\end{itemize}

21.


\begin{itemize}[leftmargin=2em,itemsep=.2em,topsep=.4em]\item[A.] It is free for us to download the app.
\item[B.] It helps you to be a professional cook.
\item[C.] It provides advice about making recipes.
\item[D.] It only rates recipes by popularity.
\end{itemize}

22.


\begin{itemize}[leftmargin=2em,itemsep=.2em,topsep=.4em]\item[A.] By showing the weight of 200 kinds of food.
\item[B.] By providing the price of 200 calories of food.
\item[C.] By picturing the food of 200 calories with weights.
\item[D.] By telling people 200 kinds of healthy food.
\end{itemize}

\textbf{Questions 23 to 25 are based on the recording you have just heard.}


23.


\begin{itemize}[leftmargin=2em,itemsep=.2em,topsep=.4em]\item[A.] About 43 percent of American adults.
\item[B.] About 18 percent of the whole population.
\item[C.] About 40 million American adults.
\item[D.] About a half million people in America.
\end{itemize}

24.


\begin{itemize}[leftmargin=2em,itemsep=.2em,topsep=.4em]\item[A.] To set a series of bans on public smoking.
\item[B.] To set the price of cigarettes properly.
\item[C.] To package the cigarettes with tips of warning.
\item[D.] To reduce the production and supply of cigarettes.
\end{itemize}

25.


\begin{itemize}[leftmargin=2em,itemsep=.2em,topsep=.4em]\item[A.] The office of the Surgeon General.
\item[B.] The Food and Drug Administration.
\item[C.] The Center for Tobacco Products.
\item[D.] The Center for Disease Control and Prevention.
\end{itemize}

\subsection*{\textbf{Part III Reading Comprehension (40 minutes)}}


\subsection*{\textbf{Section A}}


\textbf{Directions:} \textit{In this section, there is a passage with ten blanks. You are required to select one word for each blank} \textit{from a list of choices given in a word bank following the passage. Read the passage through carefully before making} \textit{your choices. Each choice in the bank is identified by a letter. Please mark the corresponding letter for each item on} \textit{Answer Sheet 2 with a single line through the center. You may not use any of the words in the bank more than once.}


Let’s say you love roller-skating. Just the thought of\uline{ 26 }on your roller-skates brings a smile to your face. You also know that roller-skating is excellent exercise. You have a \uline{27 }attitude toward it.


This description of roller-skating \uline{28 }the three components of an attitude: affect, cognition, and behavior. You love the activity; it’s great fun. These feelings \uline{29 }the affective or emotional component; they are an important ingredient in attitudes. The knowledge we have about the object constitutes the cognitive component of an attitude. You understand the health \uline{30} that the activity can bring. Finally, attitudes have a behavioral component. Our attitudes \uline{31 }us to go outside to enjoy roller-skating.


Now, we don’t want to leave you with the\uline{ 32 }that these three components always work together \uline{33} . They don’t; sometimes they clash. For example, let’s say you love pizza (affective component); however, you have high cholesterol and understand (knowledge component) that eating pizza may be bad for your health. Which behavior will your attitude result in, eating pizza or \uline{34} it? The answer depends on which component happens to be stronger. If you are walking past a pizza restaurant at lunchtime, your emotions and feelings probably will be stronger than your knowledge that pizza may not be the best food for your health. In that instance, you have pizza for lunch. If you are at home trying to decide where to go for dinner, however, the knowledge component may \uline{35} , and you decide to go where you can eat a healthier meal.


\begin{itemize}[leftmargin=2em,itemsep=.2em,topsep=.4em]\item[A.] avoiding
\item[B.] benefits
\item[C.] highlight
\item[D.] illustrates
\item[E.] impression
\item[F.] improves
\item[G.] inquiring
\item[H.] perfectly
\item[I.] positive
\item[J.] prevail
\item[K.] primarily
\item[L.] prompt
\item[M.] specifications
\item[N.] strapping
\item[O.] typical
\end{itemize}

\subsection*{\textbf{Section B}}


\textbf{Directions:} \textbf{\textit{I}}\textit{n this section, you are going to read a passage with ten statements attached to it. Each statement} \textit{contains information given in one of the paragraphs. Identify the paragraph from which the information is derived.} \textit{You may choose a paragraph more than once. Each paragraph is marked with a letter. Answer the questions by} \textit{marking the corresponding letter on Answer Sheet 2.} \textbf{The Changing Generation}


\begin{itemize}[leftmargin=2em,itemsep=.2em,topsep=.4em]\item[A.] It turns out today’s teenagers aren’t so scary after all. Results of USA WEEKEND’s Teens \& Parents survey
\end{itemize}

reveal a generation of young people who get along well with their parents and approve of the way they’re being raised. They think of their parents with affection and respect. They speak with Mom or Dad when they have a problem. Most feel that their parents understand them, and they believe their family is the No. 1 priority in their parents’ lives. Many even think their parents are cool! Although more than a third have an object in their rooms they would like to keep secret from their parents, rarely is it anything more alarming than a diary or off-color ( 低俗的) book or CD.


\begin{itemize}[leftmargin=2em,itemsep=.2em,topsep=.4em]\item[B.] Such results may seem surprising against the background of shocking incidents that color the way the mass media
\end{itemize}

portray the young. In October 2000, the same month the survey was taken, the Washington-based Center for Media and Public Affairs wrote in its publication Media Monitor that, in a recent month of TV news coverage of American youth, just 2\% of teens were shown at home, and just 1\% were portrayed in a work setting. In contrast, the criminal justice system accounted for nearly one out of every five visual backgrounds. No wonder parents worry their own kids might spin out of control once they hit the turbulent waters of adolescence.


\begin{itemize}[leftmargin=2em,itemsep=.2em,topsep=.4em]\item[C.] The overall facts ought to reassure us. The survey shows us that today’s teens are affectionate, sensible and far
\end{itemize}

happier than the angry and tortured souls that have been painted for us by stereotypes. From other sources, we also know teenage crime, drug abuse and premarital sex are in general decline. We, of course, need to pay attention to youngsters who are filled with discontent and hostility, but we should not allow these extreme cases to distort our view of most young people.


\begin{itemize}[leftmargin=2em,itemsep=.2em,topsep=.4em]\item[D.] My own research at the Stanford Center on Adolescence uses in-depth interviews with small samples of
\end{itemize}

youngsters rather than large-scale survey. Still, in my studies and others I have read, I find the same patterns as in USA WEEKEND’S survey. Today’s teenagers admire their parents and welcome parental guidance about important matters such as career choice—though certainly not Mom and Dad’s advice on matters of personal taste, such as music or fashion. When we ask teens to choose a hero, they usually select an older family member rather than a remote public figure. Most teens say they enjoy the company of both parents and friends.


\begin{itemize}[leftmargin=2em,itemsep=.2em,topsep=.4em]\item[E.] Contrary to some stereotypes, most adolescents believe they must be tolerant of differences among individuals
\end{itemize}

(though they do not always find this easy in the cliquish ( 拉帮结派的) environment of high school). Many of them volunteer for community service with disadvantaged people. One prevalent quality we have found in teens’ statements about themselves, their friends and their families is a strikingly positive emotional tone. By and large, these are very nice kids, and as the band The Who used to sing, “The kids are alright.”


\begin{itemize}[leftmargin=2em,itemsep=.2em,topsep=.4em]\item[F.] How much is today’s spirit of harmony a change from our more turbulent past? A mere generation ago, parent-
\end{itemize}

child relations were described as “the generation gap.” Yet even then reports of widespread youth rebellion were overdone: Most kids in the ’60s and ’70s shared their parents’ basic values. Still, it is true that American families are growing closer at the dawn of this new millennium ( 千年 ). Perhaps there is less to fight about, with the country in a period of tranquility and the dangers of drug abuse and other unwholesome behavior well known. Perhaps in the face of impersonal and intimidating globalization, a young person’s family feels more like a friendly haven than an oppressive trap. And perhaps parents are acting more like parents than in the recent past. Within just the past five years. I have noticed parents returning to a belief that teenagers need the guidance of elders rather than the liberal, “anything goes” mode of child-rearing that became popular in the second half of the 20th century.


\begin{itemize}[leftmargin=2em,itemsep=.2em,topsep=.4em]\item[G.] But missing from all these data is the sense that today’s young care very much about their country, about the
\end{itemize}

broader civic and political environment, or about the future of their society. They seem to be turning inward— generally in a pro-social manner, certainly with positive benefits for intimate relationships, but too often at the expense of a connection with the present and future world beyond, including the society they will one day inherit.


\begin{itemize}[leftmargin=2em,itemsep=.2em,topsep=.4em]\item[H.] Recently, we examined more than 400 essays on the “laws of life” that teens from two communities had written
\end{itemize}

as part of an educational program initiated by the John Templeton Foundation in Radnor, Pa. In those essays, and in follow-up interviews with a few of the teenagers, we found lots of insight, positive feeling and inspirational thinking. But we also found little interest in civic life beyond the tight circles of their family and immediate friends.


\begin{itemize}[leftmargin=2em,itemsep=.2em,topsep=.4em]\item[I.] For example, only one boy said he would like to be president when he grows up. When I was in high school,
\end{itemize}

dozens in my class alone would have answered differently. In fact, other recent studies have found there has never been a time in American history when so small a proportion of young people have sought or accepted leadership roles in local civic organizations. It is also troubling that voting rates among our youngest eligible voters—18- to 24-year-olds—are way down: Little more than one in four now go to the polls, even in national elections, compared with almost twice that many when 18-year-olds were first given the vote.


\begin{itemize}[leftmargin=2em,itemsep=.2em,topsep=.4em]\item[J.] In our interviews, many students viewed politics with suspicion and distaste. “Most politicians are kind of crooked
\end{itemize}

( 不诚实的),” one student declared. Another, discussing national politics, said, “I feel like one person can’t do that much, and I get the impression most people don’t think a group of people can do that much.” Asked what they would like to change in the world, the students mentioned only personal concerns such as slowing down the pace of life, gaining good friends, becoming more spiritual, becoming either more materially successful or less materially oriented (depending on the student’s values), and being more respectful of the Earth, animals and other people. One boy said, “I’d rather be concentrating on artistic efforts than saving the world or something.”


\begin{itemize}[leftmargin=2em,itemsep=.2em,topsep=.4em]\item[K.] It is fine and healthy for teens to cultivate their personal interests, and it is good news when young people enjoy
\end{itemize}

harmonious relations with their family and friends. But there is also a place in a young life for noble purposes that include a dedication to the broader society, a love of country and an aspiration to make their own leadership contributions.


\begin{itemize}[leftmargin=2em,itemsep=.2em,topsep=.4em]\item[L.] In the past, the young have eagerly participated in national service and civic affairs, often with lots of energy and
\end{itemize}

idealism. If this is not happening today, we should ask why. Our society needs the full participation of its younger citizens if it is to continue to thrive. We know the promise is there—this is a well-grounded, talented, warm-hearted group of youngsters. We have everything to gain by encouraging them to explore the world beyond their immediate experience and to prepare themselves for their turn at shaping that world.


36. Not many young people eligible for voting are interested in local or national elections these days.


37. Parents are concerned that their children may get involved in criminal offences once they reach their teens.


38. Even during the turbulent years of last century, youth rebellion was often exaggerated in the media.


39. Teenagers of today often turn to their parents for advice on such important matters as career choice.


40. The incidence of teenage crime and misbehavior is decreasing nowadays.


41. Young people should have lofty ideals in life and strive to be leaders.


42. Some young people like to keep something to themselves and don’t want their parents to know about it.


43. It is beneficial to encourage young people to explore the broader world and get ready to make it a better place.


44. Many teenagers now offer to render service to the needy.


45. Interviews with students find many of them are only concerned about personal matters.


\subsection*{\textbf{Section C}}


\textbf{Directions:} \textit{There are two passages in this section. Each passage is followed by some questions or unfinished} \textit{statements. For each of them there are four choices marked A), B), C)and D).You should decide on the best choice} \textit{and mark the corresponding letter on Answer Sheet 2 with a single line through the center.}


\subsection*{\textbf{Passage One}}


\textbf{Questions 46 to 50 are based on the following passage.}


Manufactures of products that claim to be environmentally friendly will face tighter rules on how they are advertised to consumers under changes proposed by the Federal Trade Commission.


The commission’s revised “Green Guides” warn marketers against using labels that make broad claims, like “eco- friendly.” Marketers must qualify their claims on the product packaging and limit them to a specific benefit, such as how much of the product is recycled.


“This is really about trying to cut through the confusion that consumers have when they are buying a product and that businesses have when they are selling a product,” said Jon Leibowitz, chairman of the commission.


The revisions come at a time when green marketing is on the rise. According to a new study, the number of 5 advertisements with green messages in mainstream magazines has risen since 1987, and peaked in 2008 at 10.4\%. In 2009, the number dropped to 9\%.


But while the number of advertisements may have dipped, there has been a rapid spread of eco-labeling. There are both good and bad players in the eco-labeling game.


In the last five years or so, there has been an explosion of green claims and environmental claims. It is clear that consumers don’t always know what they are getting.


A handful of lawsuits have been filed in recent years against companies accused of using misleading environmental labels. In 2008 and 2009, class-action lawsuits ( 集体诉讼) were filed against SC Johnson for using “Greenlist” labels on its cleaning products. The lawsuits said that the label was misleading because it gave the impression that the products had been certified by a third party when the certification was the company’s own.


“We are very proud of our accomplishments under the Greenlist system and we believe that we will prevail in these cases,” Christopher Beard, director of public affairs for SC Johnson, said, while acknowledging that “this has been an area that is difficult to navigate.”


Companies have also taken it upon themselves to contest each other’s green claims.


David Mallen, associate director of the Council of Better Business Bureau, said in the last two years the organization had seen an increase in the number of claims companies were bringing against each other for false or misleading environmental product claims.


“About once a week, I have a client that will bring up a new certification I’ve never even heard of and I’m in this industry,” said Kevin Wilhelm, chief executive officer of Sustainable Business Consulting. “It’s kind of a Wild West, anybody can claim themselves to be green.” Mr. Wilhelm said the excess of labels made it difficult for businesses and consumers to know which labels they should pay attention to.


46. What do the revised “Green Guides” require businesses do?


\begin{itemize}[leftmargin=2em,itemsep=.2em,topsep=.4em]\item[A.] Manufacture as many green products as possible.
\item[B.] Indicate whether their products are recyclable.
\item[C.] Specify in what way their products are green.
\item[D.] Attach green labels to all of their products.
\end{itemize}

47. What does the author say about consumers facing an explosion of green claims?


\begin{itemize}[leftmargin=2em,itemsep=.2em,topsep=.4em]\item[A.] They can easily see through the businesses’ tricks.
\item[B.] They have to spend lots of time choosing products.
\item[C.] They have doubt about current green certification.
\item[D.] They are not clear which products are truly green.
\end{itemize}

48. What was SC Johnson accused of in the class-action lawsuits?


\begin{itemize}[leftmargin=2em,itemsep=.2em,topsep=.4em]\item[A.] It gave consumers the impression that all its products were truly green.
\item[B.] It gave a third party the authority to label its products as environmentally friendly.
\item[C.] It misled consumers to believe that its products had been certified by a third party.
\item[D.] It sold cleaning products that were not included in the official “Greenlist.”
\end{itemize}

49. How did Christopher Beard defend his company’s labeling practice?


\begin{itemize}[leftmargin=2em,itemsep=.2em,topsep=.4em]\item[A.] There were no clear guidelines concerning green labeling.
\item[B.] His company’s products had been well received by the public.
\item[C.] It was in conformity to the prevailing practice in the market.
\item[D.] No law required the involvement of a third party in certification.
\end{itemize}

50. What does Kevin Wilhelm imply by saying “It’s kind of a Wild West” (Line 3, Para. 11)?


\begin{itemize}[leftmargin=2em,itemsep=.2em,topsep=.4em]\item[A.] Businesses compete to produce green products.
\item[B.] Each business acts its own way in green labeling.
\item[C.] Consumers grow wild with products labeled green.
\item[D.] Anything produced in the West can be labeled green.
\end{itemize}

\subsection*{\textbf{Passage Two}}


\textbf{Questions 51 to 55 are based on the following passage.}


America’s education system has become less a ladder of opportunity than a structure to transmit inequality from one generation to the next.


That’s why school reform is so critical. This is an issue of equality, opportunity and national conscience. It’s not just about education, but about poverty and justice.


It’s true that the main reason inner-city schools do poorly isn’t teachers’ unions, but poverty. Southern states without strong teachers’ unions have schools at least as awful as those in union states. Some Chicago teachers seem to think that they shouldn’t be held accountable until poverty is solved. There’re steps we can take that would make some difference, and Mayor Rahm Emanuel is trying some of them—yet the union is resisting.


I’d be sympathetic if the union focused solely on higher compensation. Teachers need to be much better paid to attract the best college graduates to the nation’s worst schools. But, instead, the Chicago union seems to be using its political capital primarily to protect weak performers.


There’s solid evidence that there are huge differences in the effectiveness of teachers. The gold standard study by Harvard and Columbia University scholars found that even in high-poverty schools, teachers consistently had a huge positive or negative impact.


Get a bottom 1\% teacher, and the effect is the same as if a child misses 40\% of the school year. Get a teacher from the top 20\%, and it’s as if a child has gone to school for an extra month or two. The study found that strong teachers in the fourth through eighth grades raised the skills of their students in ways that would last for decades. Just having a strong teacher for one elementary year left pupils a bit less likely to become mothers as teenagers, a bit more likely to go to college and earning more money at age 28.


How does one figure out who is a weak teacher? Yes, that’s a challenge. But researchers are improving systems to measure a teacher’s performance throughout the year, and, with three years of data, it’s usually possible to tell which teachers are failing.


Unfortunately, the union in Chicago is insisting that teachers who are laid off—often for being ineffective—should get priority in new hiring. That’s an insult to students.


Teaching is so important that it should be like other professions, with high pay and good working conditions but few job protections for bottom performers.


This isn’t a battle between garment workers and greedy bosses. The central figures in the Chicago schools strike are neither strikers nor managers but 350,000 children. Protecting the union demand sacrifices those students, in effect turning a blind eye to the injustice in the education system.


51. What do we learn about America’s education system?


\begin{itemize}[leftmargin=2em,itemsep=.2em,topsep=.4em]\item[A.] It provides a ladder of opportunity for the wealthy.
\item[B.] It contributes little to the elimination of inequality.
\item[C.] It has remained basically unchanged for generations.
\item[D.] It has brought up generations of responsible citizens.
\end{itemize}

52. What is chiefly responsible for the undesirable performance of inner-city schools?


\begin{itemize}[leftmargin=2em,itemsep=.2em,topsep=.4em]\item[A.] Unqualified teachers.
\item[B.] Lack of financial resources.
\item[C.] Unfavorable learning environment.
\item[D.] Subconscious racial discrimination.
\end{itemize}

53. What does the author think the union should do to win popular support?


\begin{itemize}[leftmargin=2em,itemsep=.2em,topsep=.4em]\item[A.] Assist the city government in reforming schools.
\item[B.] Give constructive advice to inner-city schools.
\item[C.] Demand higher pay for teachers.
\item[D.] Help teachers improve teaching.
\end{itemize}

54. What is the finding of the gold standard study by Harvard and Columbia University scholars?


\begin{itemize}[leftmargin=2em,itemsep=.2em,topsep=.4em]\item[A.] Many inner-city school teachers are not equal to their jobs.
\item[B.] A large proportion of inner-city children often miss classes.
\item[C.] Many students are dissatisfied with their teachers.
\item[D.] Students performance has a lot to do with teachers.
\end{itemize}

55. Why does the author say the Chicago union’s demand is an insult to students?


\begin{itemize}[leftmargin=2em,itemsep=.2em,topsep=.4em]\item[A.] It protects incompetent teachers at the expense of students.
\item[B.] It underestimates students’ ability to tell good teachers from poor ones.
\item[C.] It makes students feel that they are discriminated against in many ways.
\item[D.] It totally ignores students’ initiative in the learning process.
\end{itemize}

\subsection*{\textbf{Part IV Translation (30 minutes)}}


\textbf{Directions:} \textit{For this part, you are allowed 30 minutes to translate a passage from Chinese into English. You} \textit{shouldwrite your answer on Answer Sheet 2.}


深圳是中国广东省一座新开发的城市。在改革开放之前，深圳不过是一个渔村，仅有三万多人。20 世纪 80 年代，中国政府创建了深圳经济特区，作为实施社会主义市场经济的试验田。如今，深圳人口已经超过1000 万，整个城市发生了巨大的变化。 到2014 年，深圳的人均（per-capita）GDP 已达25000 美元，相当于世界上一些发达国家的水平。就综合经济实力而言，深圳居于中国顶尖城市之列。由于其独特的地位，深圳也是国内外企业家创业的理想之地。

\end{document}